problem:  
Let \( G \) be a simply connected solvable Lie group with Lie algebra \( \mathfrak{g} \), and consider the space  
\[
\mathcal{B} = \{\text{all Lie brackets on } \mathbb{R}^n \text{ making } G \text{ a solvable Lie group}\}/\mathrm{GL}(n,\mathbb{R}),
\]  
equipped with the normalized **bracket flow** evolution corresponding to Ricci flow on left-invariant metrics:
\[
\frac{d\mu}{dt} = -\pi(\mathrm{Ric}_\mu)\mu, \quad \text{rescaled so that } \|\mu\|=1 \text{ for all } t.
\]
For nilpotent Lie groups, it is known that this normalized flow always converges (modulo automorphisms) to a unique nilsoliton bracket when one exists.  

**Open problem:**  
For general **solvable but non-nilpotent** Lie groups:
1. Does there exist an open (or dense) set of initial brackets in \( \mathcal{B} \) whose normalized bracket flow trajectories converge, modulo automorphisms and scaling, to solvsoliton brackets?  
2. Does there exist a solvable Lie algebra whose normalized bracket flow admits **non-convergent recurrent trajectories** (e.g., periodic orbits or more exotic omega-limit sets) in \( \mathcal{B} \)?  
3. Is there a global Lyapunov function for the normalized flow (such as an extension of Lauret’s moment-map functional) ensuring monotone convergence, or can this monotonicity fail for certain solvable structures?

Clarifying any of these points would illuminate whether Ricci flow acts as a **canonical dynamical selector** for solvsolitons or whether fundamentally new asymptotic behaviors can occur in the non-nilpotent regime.  

Why is it a "good" problem:  
This problem advances beyond the well-understood nilpotent setting into the largely unexplored dynamics of solvable Lie groups under homogeneous Ricci flow. It precisely raises the question of whether the harmonic-analytic framework that organizes the nilpotent case—moment maps, variational principles, and soliton attractors—extends or breaks down in the solvable context. The issue links **geometric analysis**, **dynamical systems**, and **Lie algebraic geometry**, asking whether a PDE evolution can produce nontrivial recurrent behavior on an algebraic variety.  
The question is simple to state yet probes an interface where geometric flow theory meets algebraic dynamics, with possible outcomes (convergence, periodicity, chaos) that could redefine our understanding of homogeneous Ricci flow. This tight interplay of simplicity, depth, and interdisciplinarity makes it a genuinely “good” research problem.